%% 
%% Copyright 2007-2025 Elsevier Ltd
%% 
%% This file is part of the 'Elsarticle Bundle'.
%% ---------------------------------------------
\documentclass[final,3p,times]{elsarticle}

%% The amssymb package provides various useful mathematical symbols
\usepackage{amssymb}
%% The amsmath package provides various useful equation environments.
\usepackage{amsmath}
\usepackage{float}
\usepackage{subcaption}
\usepackage{tabularx}
\usepackage{booktabs} % Necesario para las lineas de las tablas
\journal{Computers in Biology and Medicine}

% 1. Paquete necesario para los colores
\usepackage[dvipsnames]{xcolor} 
\usepackage{hyperref}
\usepackage{xurl}
\usepackage{graphicx}

% Paquete para dar formato al código fuente
\usepackage{listings}
\renewcommand{\lstlistingname}{Código}
\definecolor{codegreen}{rgb}{0,0.6,0}
\definecolor{codegray}{rgb}{0.5,0.5,0.5}
\definecolor{codepurple}{rgb}{0.58,0,0.82}
\definecolor{backcolour}{rgb}{0.95,0.95,0.92}

\lstdefinestyle{mystyle}{
    backgroundcolor=\color{backcolour},   
    commentstyle=\color{codegreen},
    keywordstyle=\color{magenta},
    numberstyle=\tiny\color{codegray},
    stringstyle=\color{codepurple},
    basicstyle=\ttfamily\footnotesize,
    breakatwhitespace=false,         
    breaklines=true,                 
    captionpos=b,                    
    keepspaces=true,                 
    numbers=left,                    
    numbersep=5pt,                  
    showspaces=false,                
    showstringspaces=false,
    showtabs=false,                  
    tabsize=2
}
\lstset{style=mystyle}

% 2. Define aquí los colores que usas en tus gráficos
\definecolor{NCR}{HTML}{0000FF}
\definecolor{Edema}{HTML}{00FF00}   
\definecolor{ET}{HTML}{FF0000}     

\begin{document}

\begin{frontmatter}

%% Title, authors and addresses
\title{Proyecto Final: Sistema de integración en aplicaciones móviles para la optimización de interfaces mediante análisis de interacción}

\author[a]{Luis Carlos Zepeda Avila}

\affiliation[a]{organization={Universidad Autónoma de Baja California},
            city={Tijuana},
            state={Baja California},
            country={México}}

\end{frontmatter}

    \item \textbf{Objetivos Específicos:}
    \begin{itemize}
        \item El sistema mostrará cuántas interacciones realizo con los apartados interactivos.
        \item Mediante la entrada del usuario (clic), el sistema guardará la información en registros acumulables de clics y, una vez a la semana, enviaré dichos registros para su posterior análisis.
        \item Optimizar la experiencia de usuario (UX) de aplicaciones móviles, reduciendo el número promedio de interacciones (clics) requeridas para completar tareas clave.
    \end{itemize}
\end{itemize}

\vspace{0.5cm}
\noindent\textbf{Justificación:} \\
El principal desafío en el desarrollo de aplicaciones móviles es el espacio. Si no poseo experiencia o recursos en el diseño de interfaces, el riesgo de crear una interfaz ``sucia'' (poco clara o desordenada) para el usuario se incrementa notablemente. Para abordar esto, ofreceré una alternativa más eficiente para el análisis de la interfaz. Esto me permitirá modificarla con un objetivo concreto: reducir la cantidad de clics que los usuarios deben realizar en la aplicación.

\vspace{0.5cm}
\noindent\textbf{Hipótesis:} \\
El objetivo de implementar un sistema de recolección de clics es analizar y optimizar las interfaces de usuario. Esto me permitirá modificar las interfaces gráficas en aplicaciones móviles, logrando que las tareas clave requieran un 15\% menos de clics.

\vspace{0.5cm}
\noindent\textbf{Alcance y Límites del Sistema:}
\begin{itemize}
    \item El sistema estará disponible para múltiples desarrolladores.
    \item El sistema solo estará disponible para aplicaciones Android.
    \item El sistema será gratuito para proyectos manejados por un personal menor a cinco personas.
    \item El sistema no deberá ser visible para el usuario común ni intervenir en su uso.
    \item Por motivos de seguridad y privacidad, preguntaré al usuario al inicio de sesión si está de acuerdo con la recopilación de dichos datos; se podrá configurar manualmente en caso de cambiar de opinión.
    \item El sistema no tomará decisiones propias para la modificación de la interfaz.
    \item El sistema no sugerirá modificaciones de la interfaz directamente.
\end{itemize}

\vspace{0.5cm}
\noindent\textbf{Selección de Tecnología:}
\begin{itemize}
    \item \textbf{Lado del cliente:} Utilizaré Kotlin como lenguaje de programación. Para red, evalué Retrofit o Ktor Cliente para el envío de eventos. No rastrearé solo el ``clic'', capturaré el View ID, Timestamp, Activity/Fragment y las coordenadas X,Y.
    \item \textbf{Ingesta y backend (Planeación inicial):} Go o Node.js. Message Broker: Apache Kafka o RabbitMQ.
    \item \textbf{Almacenamiento y Análisis (Planeación inicial):} Base de datos OLAP: ClickHouse. Visualización: Grafana o Apache Superset. \textit{(Nota: Estas tecnologías fueron ajustadas en la Etapa 3 por viabilidad técnica).}
\end{itemize}

\section{Etapa 2: Definición Metodológica}

\subsection{Parte 1: Descripción del Sistema}

El sistema opera como un ecosistema donde la información fluye desde el dispositivo del usuario final hasta un panel de toma de decisiones.

\begin{itemize}
    \item \textbf{Procesos Lineales:} El flujo de Ingesta de Datos. El evento nace en el clic del usuario, pasa por el SDK (Kotlin), se envía vía API, llega al Backend y se deposita en la base de datos. Cada paso depende estrictamente del anterior.
    \item \textbf{Procesos Paralelos:} La Recolección Multiusuario. Múltiples dispositivos capturan y envían eventos de forma simultánea e independiente. Asimismo, mientras el sistema ingesta nuevos datos, el servidor de visualización puede estar procesando consultas analíticas en paralelo.
    \item \textbf{Procesos Iterativos:} El Ciclo de Optimización UX. Una vez que analizo los resultados, propongo cambios en la interfaz, los implemento y el proceso de medición vuelve a comenzar para validar si la métrica de clics mejoró.
    \item \textbf{Anidamiento:} Existe anidamiento en el proceso de Transmisión de Datos, ya que dentro de este ejecuto sub-procedimientos de autenticación y validación de conexión antes de liberar el paquete de datos.
\end{itemize}

\vspace{0.5cm}
\noindent\textbf{Representación Gráfica y Jerárquica}

\begin{itemize}
    \item \textbf{Sistema:} Sistema de Optimización de Interfaces Móviles (SOIM).
    \item \textbf{Procesos:} Recolección y captura de Eventos, Ingesta y Distribución de Mensajes, Almacenamiento y Análisis de Datos, Visualización de Mapas de Interacción.
    \item \textbf{Procedimientos:} Registro de coordenadas X,Y; Batching; Autenticación del SDK; Query de agregación.
    \item \textbf{Métodos:} Algoritmo de densidad de puntos (para Heatmaps); Protocolo de comunicación asíncrona.
    \item \textbf{Modelo Final:} Interfaz Móvil Optimizada (Propuesta de Rediseño basada en datos).
\end{itemize}

\vspace{0.5cm}
\noindent\textbf{Definición de Procesos Clave}

\textbf{Proceso A: Ingesta de Eventos en Tiempo Real}
\begin{itemize}
    \item \textbf{Entradas:} Eventos de interacción (clics), View ID, Coordenadas X,Y, Timestamp.
    \item \textbf{Actividades/Tareas:} Captura del evento, empaquetado en formato JSON, envío al broker, validación de esquema.
    \item \textbf{Salidas:} Registro confirmado en la base de datos analítica.
    \item \textbf{Recursos:} Conexión a internet, SDK en Kotlin, Servidor en la nube.
    \item \textbf{Controles/Reglas:} Solo envío datos si el usuario aceptó el aviso de privacidad; envío lotes bajo conexión Wi-Fi (regla opcional).
    \item \textbf{Indicadores:} Tasa de eventos por segundo (EPS), Latencia de ingesta (ms).
\end{itemize}

\textbf{Proceso B: Generación de Mapas de Calor (Análisis)}
\begin{itemize}
    \item \textbf{Entradas:} Datos históricos de clics acumulados.
    \item \textbf{Actividades/Tareas:} Ejecución de consultas SQL agregadas, cálculo de densidad de coordenadas, renderizado de capas visuales sobre capturas de pantalla de la App.
    \item \textbf{Salidas:} Mapa de calor visual (Heatmap), Reporte de flujo de usuario.
    \item \textbf{Recursos:} Motor de análisis, Herramientas de visualización.
    \item \textbf{Controles/Reglas:} El análisis debe filtrarse por versión de la aplicación y tipo de dispositivo.
    \item \textbf{Indicadores:} Porcentaje de reducción de clics en tareas clave, Tiempo promedio de análisis de reporte.
\end{itemize}

\subsection{Parte 2: Metodología y Objetivos}

\textbf{Metodología para alcanzar los objetivos planteados:}
\begin{itemize}
    \item \textbf{Objetivo 1 (Capturar interacciones):} Desarrollaré un SDK ligero en Kotlin. Se integrará en la aplicación móvil para registrar de forma global los eventos onClick. Capturaré View ID, X, Y y Timestamp, almacenándolos en un buffer local.
    \item \textbf{Objetivo 2 (Ingesta y almacenamiento):} Implementaré una arquitectura orientada a eventos, garantizando rendimiento óptimo en la escritura de los registros.
    \item \textbf{Objetivo 3 (Análisis de UX):} Conectaré la base de datos a una herramienta de visualización para generar Heatmaps. Esto me revelará las zonas de mayor y menor interacción para iterar el diseño con evidencia empírica.
\end{itemize}

\textbf{Resultados previstos:}
\begin{itemize}
    \item Un SDK funcional en Kotlin, capaz de integrarse de manera sencilla.
    \item Un backend preparado para recibir telemetría.
    \item Un dashboard interactivo para visualizar los mapas de calor.
    \item Datos que me permitan justificar cambios en las interfaces, reduciendo el esfuerzo cognitivo y físico del usuario.
\end{itemize}

\textbf{Asignaturas de la carrera aplicadas:}
\begin{itemize}
    \item \textbf{Análisis y Diseño de Software:} Levantamiento de requerimientos, diseño de arquitectura y gestión del proyecto.
    \item \textbf{Lenguajes de Programación Orientados a Objetos:} Uso de Kotlin y Java para la lógica del sistema.
    \item \textbf{Aplicaciones Móviles:} Uso intensivo de Android Studio y el SDK de Android.
    \item \textbf{Interacción Humano-Computadora (IHC):} Fundamentos de usabilidad y métricas de UX.
\end{itemize}

\subsection{Parte 3: Cronograma de Actividades (Inicial)}

\begin{table}[H]
    \centering
    \caption{Cronograma inicial planteado en la Etapa 2}
    \begin{tabularx}{\textwidth}{llX}
        \toprule
        \textbf{Fecha} & \textbf{Día} & \textbf{Actividad} \\
        \midrule
        21 de Abr & Martes & Definición de la estructura de la base de datos (Tablas de Clics). \\
        27 de Abr & Lunes & Crear el servidor básico en Python (Flask) que reciba datos. \\
        28 de Abr & Martes & Configurar la conexión con SQLite para guardar los clics. \\
        04 de May & Lunes & Modificar la App Android para que envíe el ID del botón al servidor. \\
        05 de May & Martes & Pruebas de envío: Presionar botones en el cel y verlos en BD. \\
        11 de May & Lunes & Crear script de Python para gráfica de "Botones más usados". \\
        12 de May & Martes & Generación de mapas de calor visuales. \\
        18 de May & Lunes & Análisis de resultados y propuesta de reubicación de botones. \\
        19 de May & Martes & Redacción del reporte final y preparación de capturas de pantalla. \\
        \bottomrule
    \end{tabularx}
\end{table}

\section{Etapa 3: Desarrollo e Implementación}

\subsection{Informe de Logros vs Cronograma Inicial}
He realizado la primera implementación del código en Kotlin para gestionar la actividad principal del usuario, esto representa la primera etapa del proyecto en el cronograma original. Sin embargo, al enfrentarme a la tecnología real, identifiqué que varias propuestas iniciales estaban fuera de mi alcance temporal y técnico actual. 

\subsection{Desafíos y Ajustes al Proyecto}
Actualmente estoy descubriendo y comprendiendo las capacidades del entorno (Android Studio). Al inicio me pareció abrumador ya que no contaba con conocimiento previo exhaustivo en el desarrollo nativo de aplicaciones móviles. 

\textbf{Soluciones aplicadas:}
\begin{itemize}
    \item El desarrollo e implementación de una API externa compleja y el uso de un Message Broker fueron reemplazados por una persistencia de datos local creando una base de datos SQLite directamente en el dispositivo.
    \item Descarté el uso de Grafana; en su lugar, agregaré un ``Modo Administrador'' dentro de la misma aplicación para poder visualizar los mapas de calor superpuestos en la pantalla original desde el propio teléfono.
\end{itemize}

Aunque mis expectativas técnicas iniciales para la infraestructura bajaron, estoy convencido de que puedo implementar esta versión más sencilla y autocontenida que cumple perfectamente con el propósito central: optimizar interfaces mediante el análisis de clics.

\newpage

\subsection{Preparación del Entorno y Desarrollo del Proceso Central}
He preparado mi entorno de desarrollo utilizando Android Studio con el lenguaje Kotlin. Como prueba de concepto inicial, generé el script funcional que habilita la vista principal \textit{Edge-to-Edge}, sentando las bases para capturar las coordenadas de la pantalla completa (necesario para el Heatmap).

\begin{lstlisting}[language=Java, caption=Prueba de concepto: Configuración inicial de la Actividad (MainActivity.kt)]
package com.example.pruebasdk

import android.os.Bundle
import androidx.activity.enableEdgeToEdge
import androidx.appcompat.app.AppCompatActivity
import androidx.core.view.ViewCompat
import androidx.core.view.WindowInsetsCompat

class MainActivity : AppCompatActivity() {
    override fun onCreate(savedInstanceState: Bundle?) {
        super.onCreate(savedInstanceState)
        enableEdgeToEdge()
        setContentView(R.layout.activity_main)
        ViewCompat.setOnApplyWindowInsetsListener(findViewById(R.id.main)) { v, insets ->
            val systemBars = insets.getInsets(WindowInsetsCompat.Type.systemBars())
            v.setPadding(systemBars.left, systemBars.top, systemBars.right, systemBars.bottom)
            insets
        }
    }
}
\end{lstlisting}

\begin{figure}[H] 
    \centering
    \includegraphics[width=1.0\textwidth]{androidstudio.png}
    
    \caption{Capturado en Android Studio}
    \label{fig:captura_main_activity}
\end{figure}

\newpage

\subsection{Nuevo Cronograma de Actividades (Ajustado)}
Debido a los desafíos técnicos, hemos propuesto un nuevo cronograma enfocado en el desarrollo del SDK local (iniciando el 4 de mayo de 2026):

\begin{table}[H]
    \centering
    \caption{Cronograma ajustado para desarrollo en entorno local}
    \begin{tabularx}{\textwidth}{lX}
        \toprule
        \textbf{Día} & \textbf{Actividad} \\
        \midrule
        \textbf{Día 1} & \textbf{Configuración y Captura Base (SDK):} Creación del módulo en Kotlin e implementación del \textit{OnTouchListener} global para detectar eventos (clics). \\
        \textbf{Día 2} & \textbf{Lógica de Coordenadas:} Captura de View ID, X, Y y Timestamp. Lógica para identificar la pantalla actual. \\
        \textbf{Día 3} & \textbf{Persistencia de Datos:} Creación de base de datos SQLite interna con auto-guardado cada 5 interacciones. \\
        \textbf{Día 4} & \textbf{Desarrollo de App de Prueba:} Crear interfaz sencilla con botones dispersos para validar la captura de coordenadas. \\
        \textbf{Día 5} & \textbf{Procesamiento de Densidad:} Algoritmo simple de agrupación (clustering) leyendo los datos de SQLite local. \\
        \textbf{Día 6} & \textbf{Visualización (Heatmap Local):} Crear Modo Admin en la app que pinte círculos de colores sobre la interfaz según la densidad de clics. \\
        \textbf{Día 7} & \textbf{Pruebas Finales:} Validar que el SDK no genere latencia. Exportación a CSV de los datos recolectados. \\
        \bottomrule
    \end{tabularx}
\end{table}

\end{document}
